\section{Design}
\label{sec:design}

\rvmt has three stages: hardware collection, offline reconstruction, and
artifact verification.
The collection stage captures architectural events below the operating
system.
The reconstruction stage joins those events with executable identity and
runtime context.
The verification stage ensures that every reported result is backed by a
concrete artifact in the evidence package.

\begin{figure*}[t]
\centering
\fbox{\begin{minipage}{0.90\textwidth}
\centering
\scriptsize
\begin{tabular}{c@{\;$\rightarrow$\;}c@{\;$\rightarrow$\;}c@{\;$\rightarrow$\;}c}
\textbf{Workload} &
\textbf{Hardware trace} &
\textbf{Trace buffer} &
\textbf{Behavior reconstruction} \\
\small RISC-V Linux process &
\small retire, trap, syscall, context &
\small measured window &
\small source-labeled events \\
\end{tabular}

\vspace{0.6em}
\begin{tabular}{c@{\;+\;}c@{\;+\;}c}
\textbf{ELF metadata} &
\textbf{Runtime process map} &
\textbf{Reference logs} \\
\small SHA-256, code ranges, load bias &
\small process and loaded objects &
\small QEMU, strace, host logs \\
\end{tabular}
\end{minipage}}
\caption{\rvmt pipeline. Hardware trace records are joined with executable metadata and runtime context. Reference logs are used for comparison, not reported as hardware output.}
\label{fig:architecture}
\end{figure*}


\subsection{Hardware Trace Collection}
\label{sec:design:collection}

\emph{Challenge: How to capture the right architectural events---retire, trap,
syscall, privilege---without software interference on a RISC-V core.}

To address this challenge, we design a CVA6 trace adapter that emits event
records for retire, trap, syscall entry, syscall return, privilege
transition, branch, jump, context, pointer-byte snapshots, and markers.
Markers define the measured workload window: the begin marker clears the
trace buffer and the end marker stops recording, isolating the target
workload from boot or background activity.
We enforce strict \texttt{SRET} qualification so that user-mode return
evidence is not widened by permissive timing rules.

Our insight is that capturing events at instruction retirement observes
committed behavior that cannot be rolled back or speculative.
This places the observation point outside the target process and below
software tracing APIs such as \texttt{ptrace} or \texttt{strace}.
Because the adapter observes committed architectural state through the
hardware trace path, user-space anti-debug checks do not directly observe
the trace adapter in the way they observe software tracers. We still treat
general hardware invisibility and trace-off slowdown as evaluation claims
that require separate evidence.

\subsection{Source-Labeled Reconstruction}
\label{sec:design:labels}

\emph{Challenge: How to ensure reconstructed behavior is auditable by
labeling each field with its provenance source.}

\begin{figure}[t]
\centering
\fbox{\begin{minipage}{0.86\columnwidth}
\scriptsize
\textbf{Source labels}

\vspace{0.4em}
\begin{tabularx}{\linewidth}{p{0.24\linewidth}Y}
Hardware & event kind, PC, trap cause, marker order \\
ELF & binary hash, section range, symbol range \\
Runtime map & process identity, load object ownership \\
Reference log & expected syscall or path behavior \\
\end{tabularx}
\end{minipage}}
\caption{Each reconstructed field records where it came from.}
\label{fig:source-labels}
\end{figure}


A common failure mode in trace analysis is reporting reconstructed fields
without identifying their source, which risks overstating what the hardware
alone recovered.
We avoid this by designing every reconstructed field to carry a provenance
label that records which source contributed to it.

We define four source labels.
(1)~\emph{Hardware trace} labels fields derived directly from the CVA6
event stream, such as program counters, syscall numbers, and trap causes.
(2)~\emph{ELF metadata} labels fields derived from binary identity and
code ranges, such as section boundaries and executable SHA-256 digests.
(3)~\emph{Runtime process map} labels fields derived from load addresses
and process context, such as object base addresses and memory-range
ownership.
(4)~\emph{Reference log} labels fields derived from QEMU, \texttt{strace},
or host-control data that validate expected behavior.

Reference-derived fields serve an important validation role but are never
treated as hardware output.
This design enables auditable behavior summaries: a reader can trace any
field back to its evidence source and verify whether the reported value
was observed directly by the hardware or supplied by supporting metadata.

\subsection{Code Attribution}
\label{sec:design:attribution}

\emph{Challenge: How to connect raw program-counter events to executable
identity and process context under PIE, ASLR, \texttt{fork}, \texttt{execve},
and dynamic loading.}

Raw PC events do not directly answer security questions without connection
to executable identity.
Simple PC-to-binary mapping fails under several conditions that are
common on modern Linux systems.

\paragraph{PIE and ASLR.}
Position-independent executables load at randomized addresses.
We resolve the load bias by reading the runtime process map and adjusting
the PC relative to the reported base, enabling accurate attribution even
when the load address differs across runs.

\paragraph{Dynamic libraries.}
Shared objects are loaded at runtime into unpredictable address ranges.
We match each PC to the loaded object ranges recorded in the runtime
process map, ensuring that code executing in a shared library is
attributed to the correct object rather than the main executable.

\paragraph{Fork and exec.}
Child processes inherit and transform address spaces across \texttt{fork}
and \texttt{execve}.
We track process ownership across these transitions by monitoring
context-switch and privilege-transition events in the trace stream,
maintaining accurate attribution even when the process tree changes.

\paragraph{Stripped binaries.}
Debug symbols are absent from deployed binaries.
We attribute code using symbol-free ELF metadata, matching PCs to section
ranges and SHA-256 binary identity without requiring DWARF information or
reconstructed source-line mappings.

Our approach combines three mechanisms to resolve these cases: ELF
SHA-256 identity for exact binary matching; runtime process maps for
load bias and object ranges; and \texttt{fork} and \texttt{execve} event
tracking for process lineage.
Source-map boundaries are handled by matching board artifacts to local
build outputs when available.
